\subsection{Completeness and composition soundness of the compiler}\label{sec:dpr:security}

\begin{lemma}[Honest initial contraction]\label{lem:dpr:init-honest}
If the public claim is true,
$y_{\rm pub}=y_w$, and the prover
sends the correct tensor message
$S_0=S$, the initial contraction
check passes.
\end{lemma}
\begin{proof}
By Theorem~\ref{thm:dpr:claim-contraction},
$C_q(S)=y_w$. Substitute
$S_0=S$ and
$y_{\rm pub}=y_w$.
\end{proof}

\begin{lemma}[Honest sumcheck telescoping]\label{lem:dpr:telescoping}
If $S_0=S=\sum_{b\in I_h}H(b)$
and the prover sends the true
partial-sum polynomial $g_i=G_i$
at every round, then all endpoint
checks pass and
\[
 s_h=H(\varphi_1(z)).
\]
\end{lemma}
\begin{proof}
Induct on the number $i$ of
completed rounds. Initially
$s_0=S$ equals the full Boolean
sum by Theorem~\ref{thm:dpr:boolean-sum}.
Suppose $s_i$ equals the true
Boolean sum over the remaining
variables at challenge prefix
$z_{<i}$. The endpoint identity
of Lemma~\ref{lem:dpr:true-partial}
gives
$s_i=G_i(0)+G_i(1)=g_i(0)+g_i(1)$,
so the verifier's check passes.
It then sets
$s_{i+1}=G_i(\varphi_1(z_i))$,
which by the definition of $G_i$
is exactly the sum over the
remaining $h-i-1$ Boolean variables.
This proves the induction invariant.
At $i=h$, no variables remain,
and the sole term is
$H(\varphi_1(z_0),\ldots,
\varphi_1(z_{h-1}))$.
\end{proof}

\begin{lemma}[Honest terminal check]\label{lem:dpr:terminal-honest}
If the sumcheck is honest and
the prover sends $u=T'(z)$,
the terminal tensor equality
in Definition~\ref{def:dpr:full-protocol}
passes.
\end{lemma}
\begin{proof}
Lemma~\ref{lem:dpr:telescoping}
gives $s_h=H(\varphi_1(z))$.
Lemma~\ref{lem:dpr:terminal}
identifies this element with
$\varphi_1(T'(z))$
times the public equality factor.
Replacing $T'(z)$ by $u$
proves the terminal check.
\end{proof}

\begin{theorem}[Perfect completeness of the unmasked compiler]\label{thm:dpr:complete}
Assume the underlying large-field
evaluation protocol has perfect
completeness in the sense of
Definition~\ref{def:dpr:large-pcs}.
For an honest Boolean witness $w$
and true public evaluation claim
$y_{\rm pub}=y_w$, an honest
execution of
Definition~\ref{def:dpr:full-protocol}
is accepted with probability $1$.
\end{theorem}
\begin{proof}
The prover constructs the packed
polynomial and commits to it.
The initial contraction passes by
Lemma~\ref{lem:dpr:init-honest}.
Every sumcheck round passes by
Lemma~\ref{lem:dpr:telescoping}.
The terminal equality passes by
Lemma~\ref{lem:dpr:terminal-honest}.
The final large-field claim
$u=T'(z)$ is true, so the
delegated evaluation proof passes
by its assumed perfect
completeness. All checks pass
for every choice of verifier
challenges.
\end{proof}

\begin{definition}[Unpacking an arbitrary committed polynomial]\label{def:dpr:extract}
For any multilinear
$T'\in L[X_0,\ldots,X_{h-1}]_{\rm mult}$,
define its underlying Boolean table
\[
 \operatorname{Dec}(T')(v,j):=
 d_v(T'(j))\in\F_2.
\]
Let $w^{T'}$ denote this table.
\end{definition}

\begin{lemma}[Packing after extraction gives the original packed polynomial]\label{lem:dpr:extract-pack}
For each $j\in I_h$,
\[
 \sum_v\beta_vw^{T'}_{v,j}=T'(j).
\]
Therefore the packed multilinear
extension of $w^{T'}$ is exactly
$T'$.
\end{lemma}
\begin{proof}
The first equality is the
unique binary basis expansion
$T'(j)=\sum_v\beta_vd_v(T'(j))$.
Both $T'$ and the packed multilinear
extension of $w^{T'}$ are
multilinear polynomials over $L$
with identical values at all
Boolean vertices $j\in I_h$.
The uniqueness of multilinear
interpolation (from Chapter~1)
therefore makes them identical.
\end{proof}

\begin{lemma}[False original claim implies false tensor initial sum]\label{lem:dpr:false-initial}
Fix a multilinear packed polynomial
$T'$ and its extracted Boolean
table $w^{T'}$.
If
$y_{\rm pub}\neq
\widetilde{w^{T'}}(r_{\rm lo},r_{\rm hi})$
and the initial contraction check
$y_{\rm pub}=C_q(S_0)$ passes,
then
\[
 S_0\neq\sum_{j\in I_h}H(j).
\]
\end{lemma}
\begin{proof}
By Lemma~\ref{lem:dpr:extract-pack},
$T'$ is the exact packing of
$w^{T'}$.
Theorem~\ref{thm:dpr:boolean-sum}
identifies the true tensor sum
with $S$. Theorems
\ref{thm:dpr:column-evaluation}
and~\ref{thm:dpr:claim-contraction},
applied to the extracted table,
give
$C_q(\sum_jH(j))
=\widetilde{w^{T'}}(r_{\rm lo},r_{\rm hi})$.
If $S_0$ were that true sum,
the passing initial contraction
would force $y_{\rm pub}$
to equal the extracted evaluation,
contrary to hypothesis.
\end{proof}

\begin{lemma}[A correct delegated opening implies the true terminal tensor value]\label{lem:dpr:correct-open}
Suppose the delegated claimed
value satisfies $u=T'(z)$.
If the terminal equality
\[
 s_h=\varphi_1(u)
 \operatorname{eq}_h(\varphi_0(r_{\rm hi}),\varphi_1(z))
\]
passes, then
$s_h=H(\varphi_1(z))$.
\end{lemma}
\begin{proof}
Substitute $u=T'(z)$ into
the terminal equality and apply
Lemma~\ref{lem:dpr:terminal},
which identifies the product
with $H(\varphi_1(z))$.
\end{proof}

\begin{theorem}[False-claim soundness with a bound underlying large-field oracle]\label{thm:dpr:compose-sound}
Fix a committed multilinear
$T'$ and its extracted Boolean
table $w^{T'}$ before any
interactive challenges.
Assume the underlying
evaluation protocol satisfies
Definition~\ref{def:dpr:large-pcs}
with error
$\varepsilon_{\rm eval}$.
If
$y_{\rm pub}\neq
\widetilde{w^{T'}}(r_{\rm lo},r_{\rm hi})$,
then for every randomized causal
prover
\[
 \boxed{
 \Pr[\mathrm{ring\ switch\ verifier\ accepts}]
 \leq\min\left\{1,\frac{2h}{q}
               +\varepsilon_{\rm eval}\right\}.}
\]
\end{theorem}
\begin{proof}
Let $E$ denote the event that
the verifier accepts its
delegated large-field evaluation
proof even though $u\neq T'(z)$.
The underlying soundness
assumption gives
$\Pr[E]\leq\varepsilon_{\rm eval}$.
We now condition on the
complement of $E$ and analyze
the ring-switching checks.

If the initial contraction
fails, the verifier rejects.
Otherwise Lemma~\ref{lem:dpr:false-initial}
proves that the initial
ring-sumcheck claim is false.
On a run in $E^c$ where the
underlying verification accepts,
the final claimed opening
satisfies $u=T'(z)$.
If the terminal tensor equality
passes, Lemma~\ref{lem:dpr:correct-open}
then forces the final sumcheck
claim to be the \emph{true}
value $H(\varphi_1(z))$.
Theorem~\ref{thm:dpr:sumcheck-sound}
bounds the probability of
that event, together with all
sumcheck checks passing, by
$2h/q$.
Thus acceptance is contained
in the union of $E$ and
the false-to-true sumcheck event.
The union bound gives
$\Pr[\mathrm{accept}]
\leq\varepsilon_{\rm eval}+2h/q$.
Cap by $1$ to finish.
\end{proof}

\begin{corollary}[Ideal large-field oracle error]\label{cor:dpr:ideal-sound}
If the delegated evaluation
is perfectly sound
($\varepsilon_{\rm eval}=0$),
then a false original evaluation
claim survives the tensor
ring-switching reduction
with probability at most
$\min\{1,2h/q\}$.
\end{corollary}
\begin{proof}
Set $\varepsilon_{\rm eval}=0$
in Theorem~\ref{thm:dpr:compose-sound}.
\end{proof}

\begin{remark}[Difference from a full IOP knowledge theorem]
Theorem~\ref{thm:dpr:compose-sound}
assumes that the committed
polynomial $T'$ has already been
bound. It does not reprove the
emulator/extractor theorem
for the actual vector-oracle
model of Diamond--Posen,
nor does it account for
all implementation-level
binding and query errors.
The original authors give
their own compiler security
argument in Theorem~4.4
of~\cite{DiamondPosen2024}.
This chapter isolates
a quantitative algebraic
soundness interface for
later formalization.
\end{remark}
